O desenvolvimento do HormuzNet demonstrou a viabilidade de aplicar conceitos clássicos de sistemas distribuídos na construção de soluções de monitoramento marítimo tolerantes a falhas e de alta disponibilidade. A transição de uma topologia com \textit{broker} centralizado (FarmNode) para uma malha descentralizada e dinâmica eliminou com eficácia os pontos únicos de falha, garantindo a continuidade do monitoramento mesmo em cenários de queda em cascata de nós.

As decisões de projeto fundamentadas na teoria --- telemetria redundante via UDP/IP \textit{Multicast} \cite{rfc768}, ordenação distribuída por Relógios Lógicos de Lamport \cite{lamport1978time} e eleição de coordenador via Algoritmo Bully \cite{tanenbaum2021redes} --- foram validadas sob injeção contínua de falhas e estresse de carga. O sistema alcançou auto-recuperação automática de \textit{brokers} e contornou com eficiência a concorrência pelo despacho de drones, mantendo filas de ocorrências íntegras e consistentes em toda a malha. O modelo de concorrência de Go com \textit{goroutines} mostrou-se especialmente adequado para a paralelização massiva exigida por dezenas de conexões simultâneas.

Como trabalhos futuros, destacam-se: a integração de canais TLS na malha P2P para proteção das mensagens de coordenação; o uso de protocolos de roteamento dinâmico para viabilizar implantações em redes WAN, eliminando a dependência de IGMP do \textit{Multicast}; e a adoção do algoritmo Raft em substituição ao Bully para maior tolerância a partições de rede. Em suma, o HormuzNet estabelece uma base funcional e robusta para sistemas de monitoramento autônomos em ambientes operacionais críticos.
